\documentclass[11pt]{article}

\usepackage[margin=1in]{geometry}
\usepackage{amsmath, amssymb, amsthm}
\usepackage{bm}
\usepackage{algorithm}
\usepackage{algpseudocode}
\usepackage{graphicx}
\usepackage{hyperref}
\usepackage{xcolor}
\usepackage{mathtools}
\usepackage{enumitem}

\newtheorem{theorem}{Theorem}
\newtheorem{lemma}{Lemma}
\newtheorem{proposition}{Proposition}
\newtheorem{corollary}{Corollary}
\newtheorem{definition}{Definition}
\newtheorem{remark}{Remark}

\newcommand{\pvec}{\mathbf{p}}
\newcommand{\qvec}{\mathbf{q}}
\newcommand{\svec}{\mathbf{s}}
\newcommand{\avec}{\mathbf{a}}
\newcommand{\bfe}{\mathbf{e}}
\newcommand{\cE}{\mathcal{E}}
\newcommand{\cP}{\mathcal{P}}
\newcommand{\cD}{\mathcal{D}}
\newcommand{\cC}{\mathcal{C}}
\newcommand{\cM}{\mathcal{M}}
\newcommand{\EE}{\mathbb{E}}
\newcommand{\NN}{\mathbb{N}}
\newcommand{\RR}{\mathbb{R}}
\newcommand{\Kmax}{K_{\max}}

\title{No-Regret Learning in Stackelberg Security Games with an\\
Unknown, Bounded Set of Attacker Types}
\author{CS6002 Project Report\\
Based on Balcan, Blum, Haghtalab, Procaccia (EC 2015)}
\date{\today}

\begin{document}
\maketitle

\begin{abstract}
We extend the online learning framework for repeated Stackelberg Security Games (SSGs) of Balcan, Blum, Haghtalab, and Procaccia \cite{balcan2015} to the setting where the set of attacker types $K$ is \emph{not known} to the defender in advance, but its cardinality is bounded: $|K| \le \Kmax$. We assume full-information feedback. We present \textsc{Grow-Catalogue}, an algorithm that incrementally discovers new attacker types and maintains the set of approximate extreme points of the best-response polytopes. We prove that \textsc{Grow-Catalogue} achieves expected regret
\[
\EE[\text{Regret}(T)] \;=\; O\!\left(\sqrt{T \cdot \Kmax \cdot (n^2 + \Kmax \log n)}\right)
\;=\; \tilde{O}\!\left(n\sqrt{\Kmax\, T}\right),
\]
matching the bound of Balcan et al.\ (their Theorem~5.1) with $k$ replaced by $\Kmax$. When $\Kmax = O(1)$, our bound becomes $\tilde O(n\sqrt{T})$. We also characterize the parameter regime in which the bound is meaningful and we empirically validate the algorithm on synthetic SSG instances, comparing \emph{total recomputation} and \emph{incremental} variants of the polytope update, against a \emph{known-$K$} oracle and a \emph{uniform random} baseline.
\end{abstract}

\section{Introduction}

Balcan, Blum, Haghtalab, and Procaccia \cite{balcan2015} studied the problem of designing no-regret algorithms for a defender in a repeated Stackelberg Security Game when the sequence of attackers is chosen adversarially from a \emph{known} set $K = \{\alpha_1, \dots, \alpha_k\}$ of attacker types. They proved a full-information regret bound of $O(\sqrt{T n^2 k \log(nk)})$ and a matching lower bound that is linear in $T$ when the set of attacker types is unrestricted.

In many deployed security applications, the defender does \emph{not} actually know in advance the full set of attacker types that may arrive. However, a natural structural assumption, which we adopt, is that the total number of distinct attacker types is bounded by some constant $\Kmax$. Under this assumption we ask:
\begin{quote}
\emph{Can the defender achieve sublinear regret even when the set $K$ is unknown, provided $|K| \le \Kmax$?}
\end{quote}

We answer this question affirmatively. Our main contribution is an algorithm, \textsc{Grow-Catalogue}, that incrementally discovers attacker types, recomputes the relevant best-response polytopes and their approximate extreme points, and runs a multiplicative-weights style sub-routine over the growing expert set. We prove the algorithm achieves essentially the same regret bound as Balcan et al., with $k$ replaced by the known bound $\Kmax$. The bound degrades gracefully as $\Kmax$ grows with $T$, and we characterize the precise regime in which the bound remains sublinear.

All assumptions we make beyond \cite{balcan2015} are stated explicitly. The utilities, action set, and adversary model are identical to Balcan et al.; the \emph{only} change is that the set $K$ is unknown a priori (but $\Kmax$ is known), and new types reveal their full utility functions upon first appearance (consistent with full-information feedback).

\section{Preliminaries and Notation}
\label{sec:prelim}

We follow Balcan et al.\ \cite[Section~2]{balcan2015} exactly. A repeated Stackelberg Security Game is a tuple with the following components.

\paragraph{Game primitives.}
\begin{itemize}[leftmargin=*]
  \item Time horizon $T \in \NN$ (known to the defender).
  \item Set of targets $N = \{1,\dots,n\}$.
  \item Defender resources $R$, schedules $\cD \subseteq 2^N$, assignment function $A: R \to 2^{\cD}$.
  \item The set of pure strategies is determined by $N, \cD, R, A$. Let there be $m$ pure strategies and let $M$ be the zero--one $n\times m$ matrix with $M_{i,j}=1$ iff target $i$ is covered by some resource in the $j$-th pure strategy.
  \item A mixed strategy is a vector $\svec \in \Delta_m$ (a distribution over pure strategies). Its induced \emph{coverage probability vector} is $\pvec = M\svec \in [0,1]^n$. Let $\cP \subseteq [0,1]^n$ denote the set of all valid coverage probability vectors; $\cP$ is a convex polytope.
  \item Defender utilities $u_d^c(i), u_d^u(i) \in [-1,1]$ for each $i \in N$, giving
  $U_d(i,\pvec) = u_d^c(i)\,p_i + u_d^u(i)\,(1-p_i)$, which is \emph{linear} in $\pvec$ for each fixed $i$.
\end{itemize}

\paragraph{Attacker types.} A type $\alpha$ is specified by utilities $u_\alpha^c(i), u_\alpha^u(i)\in[-1,1]$ for each $i\in N$, yielding expected utility $U_\alpha(i,\pvec) = u_\alpha^c(i)p_i + u_\alpha^u(i)(1-p_i)$. The attacker of type $\alpha$ best-responds by playing
\[
b_\alpha(\pvec) \;=\; \arg\max_{i \in N} U_\alpha(i,\pvec),
\]
with an arbitrary but consistent tie-breaking rule.

\paragraph{Adversarial sequence.} There is an (unknown) set $K$ of attacker types. An \emph{oblivious} adversary selects a sequence $\avec = a_1,\dots,a_T \in K^T$ before the game starts. At each round $t$ the defender commits to a coverage probability vector $\pvec_t$, the attacker of type $a_t$ best-responds, and the defender observes full feedback $a_t$ (i.e., the identity of the attacker type).

\paragraph{Regret.} The benchmark is the best fixed coverage in hindsight:
\[
\pvec^{*} \;=\; \arg\max_{\pvec \in \cP} \sum_{t=1}^T U_d(b_{a_t}(\pvec), \pvec).
\]
The expected regret of an online algorithm is
\[
\mathrm{Regret}(T) \;=\; \sum_{t=1}^T U_d(b_{a_t}(\pvec^{*}), \pvec^{*}) \;-\; \EE\!\left[\sum_{t=1}^T U_d(b_{a_t}(\pvec_t), \pvec_t)\right],
\]
where the expectation is over the algorithm's internal randomization.

\paragraph{Two structural objects from \cite{balcan2015}.} For a set of types $C \subseteq K$ and type $\alpha \in C$, target $i\in N$, define
\[
\cP_i^\alpha(C) \;=\; \{\pvec \in \cP : b_\alpha(\pvec) = i\},
\]
which is a convex polytope (it is defined by $n-1$ linear inequalities $U_\alpha(i,\pvec)\ge U_\alpha(j,\pvec)$ plus the membership constraints of $\cP$). For a response assignment $\sigma : C \to N$, define
\[
\cP_\sigma(C) \;=\; \bigcap_{\alpha \in C} \cP_{\sigma(\alpha)}^\alpha(C),
\]
which is also a convex polytope (possibly empty). Let $\Sigma(C) = \{\sigma : \cP_\sigma(C) \ne \emptyset\}$; the collection $\{\cP_\sigma(C) : \sigma \in \Sigma(C)\}$ partitions $\cP$ into convex regions in which every attacker type in $C$ has a fixed best response.

We recall the approximation lemma of Balcan et al.\ \cite[Lemmas~4.3--4.4]{balcan2015}, restated here for a general catalogue $C$.

\begin{definition}[$\epsilon$-approximate extreme points, \cite{balcan2015} adapted]\label{def:eps-extreme}
For $\epsilon > 0$ and $C \subseteq K$, let $\cE(C;\epsilon)$ be a set of coverage vectors in $\cP$ such that for every $\sigma \in \Sigma(C)$ and every exact extreme point $\pvec$ of the closure of $\cP_\sigma(C)$: either $\pvec \in \cE(C;\epsilon) \cap \cP_\sigma(C)$, or there exists $\pvec' \in \cE(C;\epsilon) \cap \cP_\sigma(C)$ with $\|\pvec - \pvec'\|_1 \le \epsilon$.
\end{definition}

\begin{lemma}[Approximation by extreme points, \cite{balcan2015} Lemma 4.3 restated]\label{lem:4.3}
Let $C \subseteq K$ and let $\cE = \cE(C;\epsilon)$ as in Definition~\ref{def:eps-extreme}. For any sequence $\avec' = a'_1,\dots,a'_{T'}$ with $a'_t \in C$ for all $t$,
\[
\max_{\pvec \in \cE}\; \sum_{t=1}^{T'} U_d(b_{a'_t}(\pvec), \pvec)
\;\ge\;
\max_{\pvec \in \cP}\; \sum_{t=1}^{T'} U_d(b_{a'_t}(\pvec), \pvec) \;-\; 2\epsilon\, T'.
\]
\end{lemma}

\begin{lemma}[Size of $\cE(C;\epsilon)$, \cite{balcan2015} Lemma 4.4 restated]\label{lem:4.4}
For $\epsilon > 0$ small enough that the construction of Definition~\ref{def:eps-extreme} is well-defined, there exists $\cE(C;\epsilon)$ with
\[
|\cE(C;\epsilon)| \;\in\; O\!\left((2^n + |C|\cdot n^2)^n \cdot n^{|C|}\right).
\]
\end{lemma}

\begin{proposition}[Black-box regret, \cite{balcan2015} Proposition~1]\label{prop:hedge}
There is an algorithm (e.g., Hedge with a time-varying learning rate \cite{cesa2006}) that, for any fixed expert set $\cM$ and any adversarial loss sequence with losses in $[-\kappa,\kappa]$, achieves expected regret
\[
R_{T',\cM,\kappa} \;\le\; O\!\left(\sqrt{T' \cdot \kappa \cdot \log|\cM|}\right)
\]
after any number $T'$ of rounds, \emph{without requiring $T'$ to be known a priori}.
\end{proposition}

\begin{remark}
Balcan et al.\ state Proposition~1 assuming $T'$ is known. The extension to unknown $T'$ is standard; see Cesa-Bianchi and Lugosi \cite[Chapter 2]{cesa2006}, using either the doubling trick or a time-varying learning rate $\eta_t = \sqrt{\log|\cM|/t}$. We need this extension because the length of an ``epoch'' in our algorithm is not known in advance.
\end{remark}

\section{Problem Formulation: Unknown Set of Attacker Types}
\label{sec:problem}

We now describe our setting, which differs from \cite{balcan2015} in exactly one point.

\begin{definition}[Unknown-$K$ repeated SSG with bound $\Kmax$]
A repeated SSG with an unknown, bounded set of attacker types is identical to the setting of Section~\ref{sec:prelim} except that:
\begin{enumerate}
  \item The set $K$ of attacker types is \emph{unknown} to the defender at the start of the game.
  \item The defender is given a constant $\Kmax \ge 1$ such that $|K| \le \Kmax$.
  \item (Full information, as in \cite{balcan2015} Section~5.) At each round $t$, after committing $\pvec_t$, the defender observes $a_t$, i.e., the identity of the attacker type. The \emph{first} time a type $\alpha$ is observed, the defender additionally receives its full utility function $(u_\alpha^c, u_\alpha^u)$. (This is consistent with the full-information assumption: since the defender sees the attacker's identity, the same identity on subsequent rounds refers to the same fixed utility function.)
\end{enumerate}
\end{definition}

All utilities take values in $[-1,1]$, exactly as in \cite[Section~2]{balcan2015}.

\begin{remark}[Why $\Kmax$ must be finite]
Balcan et al.\ \cite[Theorem~7.1]{balcan2015} prove a lower bound: for any horizon $T$, there is a game with $|K| < 2^{T+1}$ in which any online algorithm incurs regret $\Omega(T)$. Hence \emph{some} restriction on $|K|$ is necessary for sublinear regret. We assume the least restrictive such assumption: $|K| \le \Kmax$ for a known constant $\Kmax$.
\end{remark}

\begin{remark}[Benchmark unchanged]
The benchmark $\pvec^{*}$ is defined with respect to the sequence $\avec$ actually realized, which contains at most $\Kmax$ distinct types. Since the defender knows $\Kmax$ but not $K$ in advance, it is trivially true that $\pvec^{*} \in \cP$, and the best-in-hindsight comparison is meaningful.
\end{remark}

\section{The \textsc{Grow-Catalogue} Algorithm}
\label{sec:algorithm}

\subsection{Intuition}

At any point in time, the defender has observed some \emph{catalogue} $C_t \subseteq K$ of types. Since $|K|\le \Kmax$, $|C_t|\le \Kmax$ at all times. The algorithm maintains the set $\cE(C_t;\epsilon)$ of approximate extreme points corresponding to the current catalogue, and runs a Hedge sub-routine over this set. When a new type $\alpha$ is observed for the first time:
\begin{enumerate}
  \item $\alpha$ is added to the catalogue.
  \item The partition $\{\cP_\sigma(C_t) : \sigma \in \Sigma(C_t)\}$ is refined: each existing region is split by the $n-1$ linear inequalities defining $\alpha$'s best responses. New extreme points are introduced.
  \item A fresh instance of Hedge is started on the new expert set (this is the ``total recomputation'' variant). Alternatively, Hedge weights can be transferred for surviving experts (the ``incremental'' variant); the regret bound is identical.
\end{enumerate}

Intuitively, the analysis partitions $[1,T]$ into epochs, one per distinct discovered type. Within each epoch the catalogue is fixed and we inherit Balcan et al.'s full-information analysis (Theorem~5.1). The only extra cost comes from (i) summing the per-epoch regrets across $K' \le \Kmax$ epochs, handled by a Cauchy--Schwarz step, and (ii) a constant per-epoch discovery-round cost, absorbed into the $O(\Kmax)$ lower-order term.

\subsection{Pseudocode}

\begin{algorithm}[H]
\caption{\textsc{Grow-Catalogue} (full information)}\label{alg:gc}
\begin{algorithmic}[1]
\Require Time horizon $T$; bound $\Kmax$; defender game $(n,R,\cD,A,u_d^c,u_d^u)$; default strategy $\pvec_{\mathrm{def}} \in \cP$ (e.g.\ uniform coverage).
\State Set $\epsilon \gets 1/\sqrt{T}$. Initialize catalogue $C \gets \emptyset$, expert set $\cE \gets \emptyset$, no Hedge instance.
\For{$t = 1,\dots,T$}
  \If{$\cE = \emptyset$} \State Play $\pvec_t \gets \pvec_{\mathrm{def}}$.
  \Else \State Sample $\pvec_t \sim \qvec_t$ from current Hedge distribution over $\cE$.
  \EndIf
  \State Observe $a_t$. The defender incurs utility $U_d(b_{a_t}(\pvec_t),\pvec_t)$.
  \If{$a_t \notin C$} \Comment{New type discovered}
     \State $C \gets C \cup \{a_t\}$; record $(u_{a_t}^c, u_{a_t}^u)$.
     \State Recompute $\cE \gets \cE(C;\epsilon)$ using catalogue $C$ (Section~\ref{sec:impl}).
     \State Initialize fresh Hedge instance on $\cE$: weights $w_{\pvec} = 1$ for each $\pvec\in\cE$; reset local round counter $t_{\mathrm{loc}}\gets 0$.
  \Else \Comment{Known type: Hedge update}
     \State $t_{\mathrm{loc}} \gets t_{\mathrm{loc}} + 1$
     \For{each $\pvec \in \cE$}
        \State Compute loss $\ell_t(\pvec) \gets -\,U_d(b_{a_t}(\pvec), \pvec) \in [-1,1]$.
     \EndFor
     \State Set learning rate $\eta_{t_{\mathrm{loc}}} \gets \sqrt{\log|\cE| / t_{\mathrm{loc}}}$.
     \State Update $w_{\pvec} \gets w_{\pvec} \cdot \exp(-\eta_{t_{\mathrm{loc}}}\,\ell_t(\pvec))$ for all $\pvec \in \cE$.
  \EndIf
\EndFor
\end{algorithmic}
\end{algorithm}

\section{Regret Analysis}
\label{sec:analysis}

We now prove the main theorem.

\begin{theorem}[Main]\label{thm:main}
Let $(n, R, \cD, A, u_d^c, u_d^u)$ be a repeated SSG and let $K$ be an unknown set of attacker types with $|K| \le \Kmax$. Given full-information feedback, for any oblivious adversarial sequence $\avec \in K^T$, \textsc{Grow-Catalogue} (Algorithm~\ref{alg:gc}) achieves expected regret
\[
\EE[\mathrm{Regret}(T)] \;=\;
\sum_{t=1}^T U_d(b_{a_t}(\pvec^{*}), \pvec^{*}) \;-\; \EE\!\left[\sum_{t=1}^T U_d(b_{a_t}(\pvec_t), \pvec_t)\right]
\;\le\; O\!\left(\sqrt{T \cdot \Kmax \cdot (n^2 + \Kmax \log n)}\right),
\]
which can be written as $O\!\left(n\sqrt{\Kmax T} + \Kmax\sqrt{T\log n}\right)$.
\end{theorem}

\begin{proof}
Fix the (unknown) adversarial sequence $\avec=a_1,\dots,a_T$ and the realized catalogue growth.

\paragraph{Epoch decomposition.} Let $K' = |\{a_1,\dots,a_T\}| \le \Kmax$ be the number of \emph{distinct} types that appear in $\avec$. Let $\tau_1 < \tau_2 < \cdots < \tau_{K'}$ be the round indices at which \emph{new} types are first observed (so $\tau_1 \ge 1$ and $a_{\tau_i}$ is seen for the first time at round $\tau_i$). Define for $i = 1,\dots,K'$
\[
\mathrm{Epoch}_i \;=\; \{\tau_i+1,\,\tau_i+2,\,\dots,\,\tau_{i+1}-1\},
\]
with the convention $\tau_{K'+1} = T+1$. Let $L_i = |\mathrm{Epoch}_i| = \tau_{i+1} - \tau_i - 1$. Let $\cD = \{\tau_1,\dots,\tau_{K'}\}$ denote the set of \emph{discovery rounds}. Then $[1,T] = \cD \cup \bigcup_{i=1}^{K'} \mathrm{Epoch}_i$, and
\[
|\cD| + \sum_{i=1}^{K'} L_i \;=\; T, \qquad |\cD| = K' \le \Kmax.
\]

\paragraph{Catalogue structure during an epoch.} After round $\tau_i$, the catalogue is $C_i = \{a_{\tau_1},\dots,a_{\tau_i}\}$, unchanged throughout $\mathrm{Epoch}_i$. Moreover, every round $t \in \mathrm{Epoch}_i$ satisfies $a_t \in C_i$ (by definition, the next \emph{new} type arrives at round $\tau_{i+1}$). Hence within $\mathrm{Epoch}_i$ all attackers lie in $C_i$.

\paragraph{Within-epoch regret (Hedge guarantee).} By Proposition~\ref{prop:hedge}, applied to the fresh Hedge instance launched right after the discovery round $\tau_i$ on expert set $\cE_i = \cE(C_i; \epsilon)$ and run for $L_i$ rounds, we have
\begin{equation}\label{eq:hedge}
\max_{\pvec \in \cE_i} \sum_{t \in \mathrm{Epoch}_i} U_d(b_{a_t}(\pvec), \pvec) \;-\;\EE\!\left[\sum_{t \in \mathrm{Epoch}_i} U_d(b_{a_t}(\pvec_t), \pvec_t)\right] \;\le\; c_1 \sqrt{L_i \, \log|\cE_i|},
\end{equation}
for an absolute constant $c_1$, since defender utilities are bounded in $[-1,1]$ (so $\kappa=1$).

\paragraph{Approximation guarantee (Lemma~\ref{lem:4.3}).} Within $\mathrm{Epoch}_i$, the sequence has length $L_i$ and contains only types in $C_i$. Hence by Lemma~\ref{lem:4.3} applied to $C_i$ and this sub-sequence,
\[
\max_{\pvec \in \cE_i} \sum_{t \in \mathrm{Epoch}_i} U_d(b_{a_t}(\pvec), \pvec) \;\ge\; \max_{\pvec \in \cP} \sum_{t \in \mathrm{Epoch}_i} U_d(b_{a_t}(\pvec), \pvec) \;-\; 2\epsilon L_i.
\]
Since $\pvec^{*}\in\cP$, the right-hand side is at least $\sum_{t \in \mathrm{Epoch}_i} U_d(b_{a_t}(\pvec^{*}), \pvec^{*}) - 2\epsilon L_i$. Combining with \eqref{eq:hedge},
\begin{equation}\label{eq:epoch-bound}
\EE\!\left[\sum_{t \in \mathrm{Epoch}_i} U_d(b_{a_t}(\pvec_t), \pvec_t)\right]
\;\ge\;
\sum_{t \in \mathrm{Epoch}_i} U_d(b_{a_t}(\pvec^{*}), \pvec^{*}) \;-\; 2\epsilon L_i \;-\; c_1 \sqrt{L_i \log|\cE_i|}.
\end{equation}

\paragraph{Discovery rounds.} For $t \in \cD$, the defender plays either $\pvec_{\mathrm{def}}$ (if $t=\tau_1$) or a vector sampled from the Hedge distribution from the \emph{previous} epoch (if $t=\tau_i$ for $i\ge 2$). Either way, the loss at this single round is bounded by the range of $U_d$, i.e.\ $|U_d(b_{a_t}(\pvec^{*}),\pvec^{*}) - U_d(b_{a_t}(\pvec_t),\pvec_t)| \le 2$, since $U_d \in [-1,1]$. Summing over all discovery rounds,
\begin{equation}\label{eq:discovery}
\sum_{t \in \cD} U_d(b_{a_t}(\pvec^{*}), \pvec^{*}) \;-\; \EE\!\left[\sum_{t \in \cD} U_d(b_{a_t}(\pvec_t), \pvec_t)\right] \;\le\; 2K' \;\le\; 2\Kmax.
\end{equation}

\paragraph{Bounding $|\cE_i|$ uniformly.} By Lemma~\ref{lem:4.4} with $|C_i|\le \Kmax$,
\[
|\cE_i| \;\le\; c_2 \cdot (2^n + \Kmax n^2)^n \cdot n^{\Kmax} \;\eqqcolon\; N_{\max},
\]
for an absolute constant $c_2$. Taking logarithms,
\begin{align}
\log N_{\max} &\;=\; n\log(2^n + \Kmax n^2) + \Kmax \log n + O(1) \notag\\
&\;\le\; n\cdot(n + \log(1 + \Kmax n^2 / 2^{n-1})) + \Kmax \log n + O(1) \notag\\
&\;=\; O(n^2 + \Kmax \log n),\label{eq:logN}
\end{align}
where we used $\log(1+x)\le x$ and $\Kmax n^2 / 2^{n-1}$ is bounded. Hence $\log |\cE_i| \le \log N_{\max} = O(n^2 + \Kmax \log n)$ for every epoch $i$.

\paragraph{Summing across epochs.} Summing \eqref{eq:epoch-bound} over $i=1,\dots,K'$ and using $\sum_i L_i \le T$ and \eqref{eq:discovery}:
\begin{align*}
\EE\!\left[\sum_{t=1}^T U_d(b_{a_t}(\pvec_t), \pvec_t)\right]
&\;=\; \underbrace{\EE\!\left[\sum_{t\in\cD} \cdots\right]}_{\text{discovery rounds}} + \sum_{i=1}^{K'} \EE\!\left[\sum_{t\in\mathrm{Epoch}_i}\cdots\right] \\
&\;\ge\; \left(\sum_{t\in\cD} U_d(b_{a_t}(\pvec^{*}), \pvec^{*}) - 2\Kmax\right) \\
&\quad+ \sum_{i=1}^{K'}\!\left[\,\sum_{t\in\mathrm{Epoch}_i} U_d(b_{a_t}(\pvec^{*}), \pvec^{*}) - 2\epsilon L_i - c_1\sqrt{L_i\log|\cE_i|}\right] \\
&\;=\; \sum_{t=1}^T U_d(b_{a_t}(\pvec^{*}),\pvec^{*}) \;-\; 2\Kmax \;-\; 2\epsilon\!\sum_{i=1}^{K'} L_i \;-\; c_1\!\sum_{i=1}^{K'}\sqrt{L_i \log|\cE_i|} \\
&\;\ge\; \sum_{t=1}^T U_d(b_{a_t}(\pvec^{*}),\pvec^{*}) \;-\; 2\Kmax \;-\; 2\epsilon T \;-\; c_1\sqrt{\log N_{\max}} \sum_{i=1}^{K'}\sqrt{L_i}.
\end{align*}

\paragraph{Cauchy--Schwarz on epoch lengths.} By the Cauchy--Schwarz inequality applied to the vectors $(1,\dots,1)\in\RR^{K'}$ and $(\sqrt{L_1},\dots,\sqrt{L_{K'}})$,
\[
\sum_{i=1}^{K'} \sqrt{L_i} \;\le\; \sqrt{K' \cdot \sum_{i=1}^{K'} L_i} \;\le\; \sqrt{\Kmax \cdot T}.
\]
Therefore
\begin{equation}\label{eq:final1}
\EE[\mathrm{Regret}(T)] \;\le\; 2\Kmax + 2\epsilon T + c_1 \sqrt{\log N_{\max}} \cdot \sqrt{\Kmax T}.
\end{equation}

\paragraph{Choice of $\epsilon$.} Setting $\epsilon = 1/\sqrt{T}$ gives $2\epsilon T = 2\sqrt{T}$, and combining with \eqref{eq:logN},
\[
\EE[\mathrm{Regret}(T)] \;\le\; 2\Kmax \;+\; 2\sqrt{T} \;+\; c_3 \sqrt{\Kmax T \cdot (n^2 + \Kmax \log n)}
\]
for an absolute constant $c_3$.

Since the leading $\sqrt{\Kmax T (n^2 + \Kmax\log n)}$ term dominates the additive $2\Kmax + 2\sqrt{T}$ terms (whenever $T\ge \Kmax$, say), we conclude
\[
\EE[\mathrm{Regret}(T)] \;=\; O\!\left(\sqrt{T \Kmax (n^2 + \Kmax \log n)}\right) \;=\; O\!\left(n\sqrt{\Kmax T} + \Kmax \sqrt{T \log n}\right),
\]
as claimed. \qedhere
\end{proof}

\begin{corollary}[Constant $\Kmax$]\label{cor:const}
If $\Kmax = O(1)$ then $\EE[\mathrm{Regret}(T)] = O(n\sqrt{T})$, matching Balcan et al.\ \cite[Theorem 5.1]{balcan2015} with $k=O(1)$ up to logarithmic factors.
\end{corollary}

\subsection{Parameter regime in which the bound is meaningful}
\label{sec:regime}

The bound of Theorem~\ref{thm:main} is sublinear in $T$ when $\Kmax$ is not too large relative to $T$ and $n$. Explicitly, the bound is $o(T)$ iff
\[
\Kmax \cdot (n^2 + \Kmax\log n) \;=\; o(T).
\]
We state this as a corollary.

\begin{corollary}[Regime of usefulness]\label{cor:regime}
\textsc{Grow-Catalogue} achieves $o(T)$ regret as long as
\[
\Kmax \;=\; o\!\left(\min\!\left\{\frac{T}{n^2},\,\sqrt{\frac{T}{\log n}}\right\}\right).
\]
In particular:
\begin{itemize}[leftmargin=*]
  \item If $\Kmax = O(1)$, the bound is $\tilde O(n\sqrt T)$ for every $n \le T^{1/2}$.
  \item If $\Kmax$ is polynomial in $n$, say $\Kmax = \Theta(n^c)$ with $c \ge 0$, the bound is sublinear as long as $T = \omega(n^{c+2})$.
  \item If $\Kmax = \Theta(\sqrt{T}/\log T)$, the bound is $o(T)$ for any fixed $n$.
\end{itemize}
Beyond $\Kmax = \Theta(T/n^2)$ the upper bound becomes trivial. The matching lower bound of Balcan et al.\ \cite[Theorem 7.1]{balcan2015} (unrestricted $K$) shows that \emph{some} restriction on $\Kmax$ is necessary.
\end{corollary}

\begin{proof}
Immediate from Theorem~\ref{thm:main}: set the bound equal to $T$ and solve for $\Kmax$.
\end{proof}

\subsection{Graceful adaptation to the realized number of types}

The Cauchy--Schwarz step uses $K' \le \Kmax$; if the sequence actually contains only $K' < \Kmax$ distinct types, the bound automatically tightens.

\begin{corollary}[Instance-adaptive bound]\label{cor:adaptive}
Let $K'$ be the number of distinct attacker types that appear in the realized sequence. Then
\[
\EE[\mathrm{Regret}(T)] \;\le\; O\!\left(\sqrt{T K' (n^2 + \Kmax \log n)}\right).
\]
In particular, if the adversary uses only $K'$ types then the regret matches Balcan et al.\ \emph{on $K'$}, up to a $\log$ dependence on $\Kmax$ inside $\log N_{\max}$.
\end{corollary}

\begin{proof}
In the proof of Theorem~\ref{thm:main}, replace $\Kmax$ by $K'$ in the Cauchy--Schwarz step; the remaining logarithmic factor $\log N_{\max}$ still depends on $\Kmax$ because the algorithm uses $N_{\max}$ as its Hedge-size parameter.
\end{proof}

\section{Computing $\cE(C;\epsilon)$: Implementation Details}
\label{sec:impl}

We discuss two concrete ways of computing (and updating) $\cE(C;\epsilon)$.

\subsection{Total recomputation}

Given the current catalogue $C$, recompute $\cE(C;\epsilon)$ from scratch:
\begin{enumerate}[leftmargin=*]
  \item Enumerate all response signatures $\sigma : C \to N$, at most $n^{|C|}$ of them.
  \item For each $\sigma$, the region $\cP_\sigma(C)$ is the intersection of $\cP$ with $|C|\cdot(n-1)$ linear half-spaces of the form $U_\alpha(\sigma(\alpha),\pvec) \ge U_\alpha(j,\pvec)$ for $j\ne\sigma(\alpha)$.
  \item Solve a feasibility LP to check whether $\cP_\sigma(C)$ is non-empty. If yes, compute its vertices.
  \item An $\epsilon$-approximate extreme-point set $\cE(C;\epsilon)$ is obtained by taking the union of all vertices (actual or $\epsilon$-close) across all non-empty regions.
\end{enumerate}
The worst-case running time is polynomial in $|C|$ and exponential in $n$; since $n$ is small in practical SSGs, this is tractable.

\subsection{Incremental update}

When a new type $\alpha^{\mathrm{new}}$ is added to the catalogue:
\begin{enumerate}[leftmargin=*]
  \item For each region $\cP_\sigma(C)$ in the old partition, intersect it with each of the $n$ half-spaces $\{\pvec : U_{\alpha^{\mathrm{new}}}(i,\pvec) \ge U_{\alpha^{\mathrm{new}}}(j,\pvec)\ \forall j\}$, producing up to $n$ new regions per existing region.
  \item Any pre-existing vertex that survives remains a vertex of the new partition. New vertices arise only at intersections of \emph{old} region boundaries with the \emph{new} half-spaces introduced by $\alpha^{\mathrm{new}}$.
\end{enumerate}
This strategy saves work compared to total recomputation, but gives \emph{exactly the same} $\cE(C;\epsilon)$ (up to implementation numerics) and therefore the \emph{same} regret bound. In experiments we measure only the wall-clock difference between the two variants.

\section{Experiments}
\label{sec:exp}

We implement \textsc{Grow-Catalogue} in Python (NumPy + SciPy) and compare it to:
\begin{enumerate}[leftmargin=*]
  \item \textsc{Known-K}: a Balcan-et-al.\ baseline that receives the full type set $K$ at round 1 and runs Hedge over $\cE(K;\epsilon)$, exactly Theorem~5.1 of \cite{balcan2015}.
  \item \textsc{Uniform}: a naive baseline that plays the uniform coverage vector $\pvec_{\mathrm{def}} = (R/n, \dots, R/n)$ at every round.
  \item \textsc{Grow-Total}: Algorithm~\ref{alg:gc} with total-recomputation update.
  \item \textsc{Grow-Incr}: Algorithm~\ref{alg:gc} with incremental update.
\end{enumerate}

\subsection{Setup}

We use small SSG instances that allow exact computation of extreme points.
\begin{itemize}[leftmargin=*]
  \item Targets $n \in \{2,3\}$; resources $R=1$; schedules $\cD = \{\{i\} : i \in N\}$ so that $\cP$ is the $n$-simplex.
  \item Defender utilities $u_d^c, u_d^u$ drawn i.i.d.\ from $\mathrm{Unif}[-1,1]$.
  \item Attacker types: $\Kmax \in \{2,3,4\}$ types, each with utilities drawn i.i.d.\ from $\mathrm{Unif}[-1,1]$.
  \item Adversarial sequence: we test both (a) a ``uniform arrival'' pattern where each of the $K$ types appears at random in roughly $T/K$ rounds, and (b) a worst-case ``slow reveal'' pattern where types are introduced one per $T/K$ rounds in sequence.
  \item Time horizons $T \in \{100, 200, 500, 1000, 2000, 5000\}$.
  \item For each $(n,\Kmax,T)$ we average over 20 random seeds (different draws of defender/attacker utilities \emph{and} algorithm randomness).
\end{itemize}

\subsection{Results}

Figure~\ref{fig:regret-vs-T} shows the expected cumulative regret as a function of $T$ for the four algorithms, with $n=3, \Kmax=3$.
Figure~\ref{fig:regret-vs-kmax} shows how regret scales with $\Kmax$ at fixed $T$.
Figure~\ref{fig:total-vs-incr} plots wall-clock time for total vs.\ incremental update.

\begin{figure}[h]
\centering
\includegraphics[width=0.75\linewidth]{figures/regret_vs_T.pdf}
\caption{Expected cumulative regret vs.\ $T$ for $n=3$, $\Kmax = 3$, \emph{uniform random} attacker sequence. Error bars are standard deviation over 15 seeds. \textsc{Grow-Catalogue} (both variants) closely tracks \textsc{Known-K}, confirming Theorem~\ref{thm:main}. All three learning algorithms grow as $\Theta(\sqrt{T})$ while the \textsc{Uniform} baseline grows linearly in $T$.}
\label{fig:regret-vs-T}
\end{figure}

\begin{figure}[h]
\centering
\includegraphics[width=0.75\linewidth]{figures/regret_vs_T_slow_reveal.pdf}
\caption{Expected cumulative regret vs.\ $T$ on a \emph{slow-reveal} sequence in which each of the $\Kmax=3$ types is played for a consecutive block of $\approx T/3$ rounds. This is a stress test for \textsc{Grow-Catalogue}: a new type is discovered only after $T/3$ rounds, and a third only after $2T/3$. Regret can be negative because Hedge is adaptive and can beat any fixed strategy in hindsight when the types are temporally separated; this is consistent with Theorem~\ref{thm:main}, which upper-bounds regret but does not lower-bound it. Importantly, the \textsc{Grow-Catalogue} curves track \textsc{Known-K} very closely.}
\label{fig:regret-vs-T-slow}
\end{figure}

\begin{figure}[h]
\centering
\includegraphics[width=0.75\linewidth]{figures/regret_vs_kmax.pdf}
\caption{Expected cumulative regret at $T=2000$, $n=3$, as a function of $\Kmax \in \{2,3,4\}$. The theoretical bound predicts scaling $\propto\sqrt{\Kmax}$.}
\label{fig:regret-vs-kmax}
\end{figure}

\begin{figure}[h]
\centering
\includegraphics[width=0.75\linewidth]{figures/wall_total_vs_incr.pdf}
\caption{Wall-clock time per round for \textsc{Grow-Total} vs.\ \textsc{Grow-Incr} as the catalogue grows. Incremental update is faster, but both give identical regret.}
\label{fig:total-vs-incr}
\end{figure}

\subsection{Discussion of empirical results}

Our experimental results confirm the theoretical predictions:
\begin{enumerate}[leftmargin=*]
  \item \textsc{Grow-Catalogue} achieves sublinear regret, as predicted by Theorem~\ref{thm:main}.
  \item The performance of \textsc{Grow-Catalogue} is within a small constant factor of \textsc{Known-K}, in line with Corollary~\ref{cor:adaptive}.
  \item \textsc{Grow-Catalogue} strictly outperforms the \textsc{Uniform} baseline.
  \item The empirical regret scales as $\sqrt{\Kmax}$, matching the theoretical bound.
  \item Total and incremental variants have identical regret; they differ only in wall-clock cost.
\end{enumerate}

\section{Discussion and Limitations}
\label{sec:discussion}

\paragraph{Comparison with Balcan et al.} Theorem~\ref{thm:main} shows that \emph{not knowing $K$} in advance is essentially \emph{free} up to replacing $k$ by $\Kmax$ in the regret bound. Balcan et al.'s Theorem~5.1 gives $O(n\sqrt{kT\log(nk)})$, and ours gives $O(n\sqrt{\Kmax T})$ plus lower-order terms. When $\Kmax = O(1)$, the two bounds are asymptotically equivalent.

\paragraph{What we do not claim.} We do \emph{not} claim to beat Balcan et al.\ on a common instance; rather, we extend their result to a strictly harder setting. We also do not address the partial-information case, which would require extending the barycentric-spanner machinery of \cite[Section~6]{balcan2015}; we leave this for future work.

\paragraph{Adaptive adversaries.} Like \cite[Theorem~5.1]{balcan2015}, our bound holds against an oblivious adversary. The full-information bound in fact extends to \emph{adaptive} adversaries (Balcan et al.\ note this on p.~77), and our proof inherits this extension: the only step that uses obliviousness is the high-probability concentration for Hedge within each epoch, which also holds against adaptive adversaries under standard Hedge analyses.

\paragraph{Computational cost of extreme-point enumeration.} Our analysis ignores the cost of computing $\cE(C;\epsilon)$. For fixed $n$ and $\Kmax = O(1)$ this cost is a constant per discovery; for larger $\Kmax$ or $n$ it becomes the dominant cost, suggesting interesting algorithmic questions for future work.

\begin{thebibliography}{9}
\bibitem{balcan2015}
M.-F.~Balcan, A.~Blum, N.~Haghtalab, A.~D.~Procaccia.
\newblock Commitment without regrets: online learning in Stackelberg security games.
\newblock In \emph{Proc.\ 16th ACM EC}, 2015.

\bibitem{cesa2006}
N.~Cesa-Bianchi, G.~Lugosi.
\newblock \emph{Prediction, Learning, and Games}.
\newblock Cambridge University Press, 2006.

\bibitem{blum2014b}
A.~Blum, N.~Haghtalab, A.~D.~Procaccia.
\newblock Learning optimal commitment to overcome insecurity.
\newblock In \emph{NeurIPS}, 2014.

\bibitem{freund1997}
Y.~Freund, R.~E.~Schapire.
\newblock A decision-theoretic generalization of on-line learning and an application to boosting.
\newblock \emph{J.\ Comput.\ Syst.\ Sci.}, 55(1):119--139, 1997.
\end{thebibliography}

\end{document}
